% ============================================================
% main.tex — selve notatet
%
% Denne fila er struktur og innhold. Alle typografiske
% valg bor i notatstil.sty, som også laster babel.
%
% Kompiler med pdfLaTeX:
%   latexmk main
% ============================================================


\documentclass[a4paper,11pt]{article}


% ============================================================
% LAST INN NOTATSTIL
% ============================================================
\usepackage{notatstil}
% Kommenter inn for å fjerne lengefarger (ren typografi, ingen farger):
% \hypersetup{hidelinks}


% ============================================================
% EGENDEFINERTE PAKKER
%
% Legg til prosjektspesifikke pakker her.
% Ikke sett pakker andre steder i filen.
% ============================================================

% --- Matematikk (standard) ---
% amsthm og teorem-/definisjon-/merknad-omgivelsene
% leveres av notatstil.sty.
\usepackage{amsmath, amssymb}
\usepackage{mathtools}

% --- Valgfrie fysikk- / matematikkutvidelser ---
% \usepackage{physics}              % \dd, \ket, \bra, \comm, ...
% \usepackage{bm}                   % uthevede mattesymboler
% \usepackage{slashed}              % Feynman-skråstrek
% \usepackage{tensor}               % indekserte tensorer

% --- Annet ---
\usepackage{graphicx}               % figurer
\graphicspath{{figurer/}}           % figurer ligger i figurer/
\usepackage{booktabs}               % tabeller
\usepackage{float}                  % [H]-plassering for figurer
\usepackage{tikz}                   % native TikZ-figurer
% \usepackage{enumitem}             % listetilpasning


% ============================================================
% EGENDEFINERTE KOMMANDOER
%
% Prosjektspesifikke makroer.
% ============================================================

% \newcommand{\dd}{\mathrm{d}}
% \newcommand{\ii}{\mathrm{i}}


% ============================================================
% METADATA
% ============================================================
\title{Notattittel}
\author{Forfatternavn}
\date{\today}


% ============================================================
% DOKUMENT
% ============================================================
\begin{document}

\maketitle

\begin{abstract}
Disse notatene utvikler en liten bit av teoretisk maskineri og
arbeider gjennom konsekvensene i detalj. Framstillingen holdes
selvstendig, med definisjoner, formuleringer, korte bevis og
gjennomgåtte eksempler vekslet med prosa. De avsluttende
seksjonene samler noen oppgaver til øving.
\end{abstract}

\clearpage
\framedtoc
\clearpage


% ============================================================
\section{Bakgrunn}

Disse notatene utvikler en liten bit av teoretisk maskineri
for å illustrere strukturen i malen. Framstillingen følger
det typiske mønsteret for matematikk- og fysikknotater:
definisjoner, formuleringer av resultater, korte bevis, og
gjennomgåtte eksempler vekslet med prosa.

Siteringer vises i hakeparenteser, f.eks.\ \cite{eksempel-artikkel},
og flere kilder kombineres og komprimeres til intervaller
automatisk~\cite{eksempel-artikkel,eksempel-bok}.


% ============================================================
\section{Definisjoner og resultater}

% viktig-boksen løfter en sentral definisjon fra prosaen;
% nummereringen fortsetter i samme rekke som utenfor.
\begin{viktig}
\begin{definisjon}
  La $\mathcal{H}$ være et komplekst separabelt Hilbertrom. En
  \emph{tilstand} er en enhetsvektor $\psi \in \mathcal{H}$,
  betraktet opp til en fasefaktor $e^{i\theta}$.
\end{definisjon}
\end{viktig}

\begin{definisjon}
  En \emph{observabel} er en selvadjungert lineær operator
  $\hat{A}$ som virker på $\mathcal{H}$.
\end{definisjon}

Forventningsverdien til en observabel i en gitt tilstand er
definert ved
\begin{equation}
  \langle A \rangle_\psi
  = \langle \psi | \hat{A} | \psi \rangle .
  \label{eq:forventning}
\end{equation}

% hovedresultat rammer inn dokumentets sentrale resultat;
% tittelargumentet kan sløyfes for en ren ramme.
\begin{hovedresultat}[Reelle forventningsverdier]
\begin{teorem}
  Forventningsverdier av selvadjungerte operatorer er reelle.
\end{teorem}
\end{hovedresultat}

\begin{proof}
  Per definisjon,
  $\overline{\langle A \rangle_\psi}
  = \overline{\langle \psi | \hat{A} | \psi \rangle}
  = \langle \psi | \hat{A}^\dagger | \psi \rangle
  = \langle \psi | \hat{A} | \psi \rangle
  = \langle A \rangle_\psi$,
  hvor vi brukte $\hat{A}^\dagger = \hat{A}$.
\end{proof}

\begin{merknad}
  At forventningsverdier er reelle, er det som gjør
  selvadjungerte operatorer til den naturlige matematiske
  modellen for observerbare størrelser i kvantemekanikk.
\end{merknad}

\subsection{Tidsutvikling}

Tidsutvikling i kvantemekanikk styres av Schr\"odinger-ligningen:
\begin{equation}
  i\hbar \frac{\partial \psi}{\partial t}
  = \hat{H}\psi ,
  \label{eq:schrodinger}
\end{equation}
hvor $\hat{H}$ er Hamilton-operatoren. For tidsuavhengig
$\hat{H}$ integrerer ligning~\eqref{eq:schrodinger} til
\begin{equation}
  \psi(t) = e^{-i\hat{H}t/\hbar}\,\psi(0).
\end{equation}

\subsubsection{Stasjonære tilstander}
En \emph{stasjonær tilstand} er en egentilstand til $\hat{H}$;
tidsavhengigheten er en ren fase, så enhver forventningsverdi
er konstant i tid. Run-in-overskrifter som denne markerer det
laveste strukturnivået og holdes utenfor innholdsfortegnelsen.


% ============================================================
\section{Eksempler}

\begin{figure}[t]
  \centering
  \includegraphics[width=0.7\linewidth]{Energy_levels}
  \caption{Skjematisk energinivådiagram.}
  \label{fig:nivaer}
\end{figure}

Figur~\ref{fig:nivaer} viser en typisk struktur som dukker
opp gjennom resten av notatene. Noen fundamentale konstanter
er samlet i tabell~\ref{tab:konstanter} for referanse.

\begin{table}[ht]
  \centering
  \caption{Fundamentale fysiske konstanter.}
  \label{tab:konstanter}
  \begin{tabular}{lcr}
    \toprule
    Størrelse           & Symbol  & Verdi                         \\
    \midrule
    Lysets hastighet    & $c$     & $2{,}998 \times 10^8$ m/s     \\
    Plancks konstant    & $\hbar$ & $1{,}055 \times 10^{-34}$ Js  \\
    Elementærladning    & $e$     & $1{,}602 \times 10^{-19}$ C   \\
    \bottomrule
  \end{tabular}
\end{table}


% ============================================================
\section{Numeriske metoder}

% TikZ-figureksempel: bruk darkorange til energinivåer / strukturelementer,
% darkolive, darkblue og darkplum til kurver / bølgefunksjoner (i den
% rekkefølgen), svart til akser og vegger, og grå til sekundære etiketter.
% darkred er reservert for hyperlenker og separatorlinjen -- bruk den
% ikke i figurer.
\begin{figure}[H]
  \centering
  \begin{tikzpicture}[>=stealth]
    % Brønnvegger og gulv
    \draw[very thick] (0,0) -- (0,4);
    \draw[very thick] (4,0) -- (4,4);
    \draw[thick]      (-0.3,0) -- (4.3,0);
    % Energiakse
    \draw[->,thick] (-0.8,0) -- (-0.8,4.2)
      node[above,font=\small] {$E$};
    % Energinivåer (darkorange)
    \foreach \n/\y in {1/0.65, 2/1.70, 3/3.10} {
      \draw[darkorange,thick] (0,\y) -- (4,\y);
      \node[right,font=\small,darkorange] at (4.15,\y) {$E_{\n}$};
    }
    % Bølgefunksjoner (darkolive, darkblue, darkplum)
    \draw[darkolive,thick,domain=0:4,samples=80]
      plot (\x, {0.65 + 0.38*sin(deg(\x*3.14159/4))});
    \draw[darkblue,thick,domain=0:4,samples=80]
      plot (\x, {1.70 + 0.38*sin(deg(2*\x*3.14159/4))});
    \draw[darkplum,thick,domain=0:4,samples=80]
      plot (\x, {3.10 + 0.38*sin(deg(3*\x*3.14159/4))});
    % x-aksemerker (grå for sekundær info)
    \node[below,font=\small,gray] at (0,-0.1) {$0$};
    \node[below,font=\small,gray] at (4,-0.1) {$L$};
  \end{tikzpicture}
  \caption{Uendelig dyp brønn: energinivåer $E_n \propto n^2$
           (\textcolor{darkorange}{oransje}) og tilhørende
           egenfunksjoner $\psi_1$, $\psi_2$ og $\psi_3$
           (\textcolor{darkolive}{olivengrønn}, \textcolor{darkblue}{blå}
           og \textcolor{darkplum}{plomme}).}
  \label{fig:bronnn}
\end{figure}

\verb|python|-omgivelsen setter kode med Gruvbox Light
syntaksutheving direkte i notatene — godt egnet til korte skript
som følger en analytisk utledning. Eksempelsystemet gjennom hele
malen er den uendelig dype brønnen, vist i figur~\ref{fig:bronnn}
med de tre første energinivåene og egenfunksjonene.

Skriptet nedenfor evaluerer sannsynlighetstettheten
$|\psi_n(x)|^2$ numerisk, og integrerer
$\langle x \rangle = \int_0^L x\,|\psi_n|^2\,\mathrm{d}x$
for de tre første nivåene og bekrefter resultatet $\langle x \rangle = L/2$.

\begin{python}
import numpy as np

def psi(n, x, L=1.0):
    """Normert bølgefunksjon for uendelig dyp brønn."""
    return np.sqrt(2.0 / L) * np.sin(n * np.pi * x / L)

def energi(n, L=1.0, hbar=1.0, m=1.0):
    """Energiegenverdi E_n i naturlige enheter."""
    return (n * np.pi * hbar)**2 / (2 * m * L**2)

x = np.linspace(0, 1, 1000)
for n in range(1, 4):
    tetthet = psi(n, x)**2          # |psi_n|^2
    middelx = np.trapz(x * tetthet, x)
    print(f"n={n}  E_n={energi(n):.4f}  <x>={middelx:.4f}")
\end{python}

Kjøring av dette gir $\langle x \rangle \approx 0{,}5000$ for
alle~$n$, i samsvar med symmetrien til hver egentilstand om
midtpunktet av brønnen.


% ============================================================
\section{Oppgaver}

Notater kan bære oppgaver ved siden av framstillingen.
\verb|oppgave|-omgivelsen nummereres per seksjon og settes i
full bredde, skilt fra teksten rundt av luft alene — mindre
før enn etter, så oppgaven bindes til innledningen sin.
\verb|\separator| kan i tillegg markere et temaskifte, som
under. Deloppgaver bruker \verb|deloppgaver|-omgivelsen, som
merker delene a), b), c) uten å påvirke vanlige
\verb|enumerate|-lister.

\separator

\begin{oppgave}[Egenverdier og ortogonalitet]
  Vis at egenverdiene til en selvadjungert operator $\hat{A}$
  er reelle, og at egenvektorer som hører til forskjellige
  egenverdier er ortogonale.
\end{oppgave}

\begin{oppgave}[Uendelig dyp brønn]
  Betrakt en partikkel i en endimensjonal uendelig dyp
  brønn med bredde $L$.
  \begin{deloppgaver}
    \item Skriv opp de normerte energiegenfunksjonene.
    \item Regn ut energinivåene $E_n$.
  \end{deloppgaver}

  De resterende delene bygger på spekteret funnet ovenfor; bruk
  resultatet fra deloppgave~(b) når du skisserer tetthetene.

  \begin{deloppgaver}[resume]
    \item Skisser sannsynlighetstettheten for $n = 1, 2, 3$.
    \item Regn ut forventningsverdien $\langle x \rangle$ i
          grunntilstanden.
  \end{deloppgaver}
\end{oppgave}


% ============================================================
\section{Avslutning}

Disse notatene er en mal; innholdet over finnes kun for å
demonstrere typografi. Bytt det ut med ditt eget materiale
etter hvert som notatene utvikler seg.


% ============================================================
% BIBLIOGRAFI
% ============================================================
\bibliographystyle{unsrtnat}
\bibliography{referanser}


\end{document}
